\section{Commutative canonicalization and duplicate handling}
\label{sec:canonicalization}

\subsection{What canonicalization means}

The tensor representation is ordered, but the scalar observables are ordinary
commutative monomials. Consequently,
\[
S_1I_1I_2=S_1I_2I_1.
\]
The two products correspond to the same scalar observable and must occupy one
coordinate in the dictionary. By contrast,
\[
S_1I_1^2
\]
contains a repeated factor and remains a distinct monomial.

A practical canonical representation stores the exponent of every scalar
variable in a fixed order. For example, with the variable order
\[
S_1,S_2,S_3,E_1,E_2,E_3,I_1,I_2,I_3,
\]
the monomial $S_1I_1I_2$ is represented by the exponent vector
\[
(1,0,0,0,0,0,1,1,0).
\]
This representation is insensitive to factor order and preserves repeated
factors exactly.

\subsection{Two distinct notions of duplication}

It is important to distinguish two cases.

\begin{enumerate}[leftmargin=2em]
\item \textbf{Duplicate terms within one scalar derivative.} If the same
      canonical monomial is produced more than once in one derivative equation,
      its coefficients must be added before the operator row is assembled.
\item \textbf{The same coordinate generated in different scalar derivatives.}
      The dictionary still contains one coordinate, but each derivative has its
      own matrix entry in the column associated with that coordinate.
\end{enumerate}

In the present degree-two-to-degree-three calculation, no exact within-row
duplicate or cancellation remains when the channel weights are retained
symbolically. Nine cubic coordinates are, however, shared by two different
scalar derivative equations:

% Generated by scripts/generate_example_artifacts.py
\begin{center}
\begin{tabular}{@{}p{0.25\textwidth}p{0.65\textwidth}@{}}
\toprule
Shared cubic coordinate & Generating scalar derivatives\\
\midrule
$S_1I_1I_2$ & $\LieQ(S_1I_1)$ and $\LieQ(S_1I_2)$ \\
$S_2I_1I_2$ & $\LieQ(S_2I_2)$ and $\LieQ(S_2I_1)$ \\
$S_2I_1I_3$ & $\LieQ(S_2I_3)$ and $\LieQ(S_2I_1)$ \\
$S_2I_2I_3$ & $\LieQ(S_2I_2)$ and $\LieQ(S_2I_3)$ \\
$S_3I_2I_3$ & $\LieQ(S_3I_3)$ and $\LieQ(S_3I_2)$ \\
$S_1S_2I_1$ & $\LieQ(S_1E_2)$ and $\LieQ(S_2E_1)$ \\
$S_1S_2I_2$ & $\LieQ(S_1E_2)$ and $\LieQ(S_2E_1)$ \\
$S_2S_3I_2$ & $\LieQ(S_2E_3)$ and $\LieQ(S_3E_2)$ \\
$S_2S_3I_3$ & $\LieQ(S_2E_3)$ and $\LieQ(S_3E_2)$ \\
\bottomrule
\end{tabular}
\end{center}


For example, the column associated with $S_1I_1I_2$ contains
\[
-\beta a_{12}
\]
in the row for $\frac{d}{dt}(S_1I_1)$ and
\[
-\beta a_{11}
\]
in the row for $\frac{d}{dt}(S_1I_2)$. These coefficients are not added
together because they belong to different equations.

After merging the nine repeated coordinates across different rows, the
\NumRawCubic{} appearances reduce to
\[
\NumRawCubic-9=\NumMthreeSeeds
\]
distinct cubic seeds. They are distributed as follows:
\[
12\text{ of type }SII,
\qquad
17\text{ of type }SIE,
\qquad
13\text{ of type }SSI.
\]
Thus
\[
|\seedM_3^{\mathcalE}|=\NumMthreeSeeds.
\]

\stepbox{Why the seed count is not yet the final cubic dimension}{
The 42 seeds are all cubic coordinates created directly by the nonlinear
transmission mechanism. They need not form an invariant degree-three block
under the linear SEIR transitions. Algorithm~1 therefore applies linear
closure before defining the final cubic dictionary.}
